\usepackage{fancyhdr}
\pagestyle{fancy}
\fancypagestyle{plain}{}
\renewcommand{\chaptermark}[1]{\markboth{#1}{}}
%\renewcommand{\sectionmark}[1]{\markright{#1}{}}

\fancyhead{}
\fancyhead[L]{
	\setlength{\unitlength}{1pt}
	\qquad \quad Universidad Tecnológica Nacional\\
	\qquad \quad Facultad Regional Reconquista\\
	\qquad \quad Año 2021\\      %para borrar el año borrar esta linea y el vspace de abajo
	\vspace*{-15pt}
	\begin{picture}(1,1)
		\put(0,-2){\includegraphics[height=39pt]{imagenes/LogoUTN_nvgsb.eps}}
\end{picture}}

\fancyhead[C]{}
\fancyhead[R]{Proyecto Final\\TU NOMBRE COMPLETO\\\leftmark}

\fancyfoot{}
\fancyfoot[L,C]{}
\fancyfoot[R]{\thepage}


\usepackage{etoolbox} % https://tex.stackexchange.com/questions/39988/space-before-chapters-and-contents
\makeatletter
\patchcmd{\@makechapterhead}{50\p@}{0pt}{}{}
\patchcmd{\@makeschapterhead}{50\p@}{0pt}{}{}
\makeatother

% Las 3 siguientes lineas no se usan, pero si se usa comando \maketitle son utiles
\title{Diseño eléctrico, programación \\ y control de servicios \\ auxiliares para caldera de \\ central térmica}
\author{Santiago Andrés Franzoi}
\date{Marzo 2021}